% Volume VII: Deep Learning Mathematics
% Continuity note for Chapter 40.
% Previous dependencies: Chapters 7 and 38 introduced tensor shapes, contractions, convolution, computational graphs, and backpropagation; Chapter 39 introduced training dynamics and gradient-flow concerns.
% New durable concepts: convolutional weight sharing, receptive fields, translation equivariance, recurrent hidden states, state-space models, attention as content-based weighted averaging, transformer blocks, graph message passing, and permutation equivariance.
% Forward hooks: Chapter 41 uses these architectures as encoders, decoders, generators, discriminators, and denoisers; Chapter 42 analyzes their learned representation geometry and robustness.
\section{Chapter 40: Deep Architectures: CNNs, RNNs, Transformers, and GNNs}
\label{ch:deep-architectures}

A neural-network architecture is a mathematical constraint on information flow. A fully connected layer treats all input coordinates as exchangeable once their positions are fixed. A convolutional network assumes locality and weight sharing. A recurrent network carries state through time. A transformer routes information by content-dependent attention. A graph neural network updates node features by messages along edges.

\paragraph{Chapter problem.}
How do modern deep architectures encode assumptions about space, time, sets, sequences, and graphs while remaining differentiable tensor programs trainable by backpropagation?

\paragraph{Section map.}
Section 40.1 studies convolutional neural networks as local, weight-sharing maps with translation equivariance. Section 40.2 studies recurrent neural networks and state-space models as learned dynamical systems over sequences. Section 40.3 derives attention and transformer blocks from query-key-value tensor operations. Section 40.4 studies graph neural networks as permutation-equivariant message-passing systems.

\subsection{40.1 Convolutional neural networks}

\paragraph{Guiding question.}
How can a neural network exploit local spatial structure without learning a separate weight for every pixel position?

\paragraph{Beginner-facing intuition.}
Images, spectrograms, and many spatial neural recordings have local structure. Nearby pixels or sensors tend to be more directly related than faraway ones. A convolutional neural network (CNN) encodes this assumption by applying the same small filter at many locations. This is called weight sharing.

The same edge detector can be useful in the top-left and bottom-right of an image. A convolutional layer therefore learns local patterns rather than a separate parameter for every absolute location. Pooling, striding, and stacked layers enlarge the receptive field so later features summarize broader regions.

\paragraph{Formal definitions and notation.}
For a one-dimensional signal \(x\in\mathbb{R}^{T\times C_{\mathrm{in}}}\), a convolutional layer with kernel radius \(r\) and \(C_{\mathrm{out}}\) output channels can be written
\[
y_{t,c_{\mathrm{out}}}
=b_{c_{\mathrm{out}}}
+\sum_{\Delta=-r}^{r}\sum_{c_{\mathrm{in}}=1}^{C_{\mathrm{in}}}
K_{\Delta,c_{\mathrm{in}},c_{\mathrm{out}}}x_{t+\Delta,c_{\mathrm{in}}},
\]
with boundary handling determined by padding. For images,
\[
X\in\mathbb{R}^{H\times W\times C_{\mathrm{in}}},
\qquad
Y\in\mathbb{R}^{H'\times W'\times C_{\mathrm{out}}},
\]
and a two-dimensional kernel uses offsets \((\Delta_h,\Delta_w)\). The \emph{stride} controls how far the filter moves between output locations. The \emph{receptive field} of an output unit is the set of input locations that can affect it.

\paragraph{Core derivation: translation equivariance.}
Let \(T_s\) shift a one-dimensional signal by \(s\) positions, so \((T_s x)_t=x_{t-s}\). Ignoring boundaries, define convolution
\[
(K\ast x)_t=\sum_{\Delta=-r}^{r}K_\Delta x_{t+\Delta}.
\]
Then
\[
(K\ast(T_s x))_t
=\sum_\Delta K_\Delta (T_s x)_{t+\Delta}
=\sum_\Delta K_\Delta x_{t+\Delta-s}
=(K\ast x)_{t-s}
=(T_s(K\ast x))_t.
\]
Thus shifting the input shifts the output by the same amount. This is \emph{translation equivariance}. It is not invariance: the feature map still moves. Pooling or global averaging can turn equivariant features into more invariant summaries.

\paragraph{Concrete example.}
Let \(x=(1,3,2,0)\) and \(K=(-1,1)\), with valid one-dimensional convolution
\[
y_t=-x_t+x_{t+1}.
\]
Then
\[
y=(2,-1,-2).
\]
The filter reports local increases and decreases. The same two weights were used at every location, so the layer used two kernel parameters rather than a separate pair of weights for each output position.

\paragraph{Computational neuroscience example.}
Simple cells in primary visual cortex are often modeled as localized oriented filters. A CNN filter is a mathematical cousin of this idea: it detects a local pattern and reuses the detector across locations. The comparison should be stated carefully. CNN weight sharing is an engineered symmetry; cortical receptive fields are not literally one learned kernel copied perfectly across the retina.

\paragraph{AI/ML example.}
In image classification, early CNN layers often learn edge-like and color-contrast filters; later layers combine them into object parts and category-relevant features. In audio or neural time-series analysis, one-dimensional convolutions can detect local temporal motifs such as bursts, oscillation fragments, or waveform shapes.

\paragraph{Common misconceptions and failure modes.}
A convolution is not automatically invariant to translations; it is equivariant before pooling or aggregation. A large CNN can still overfit if data are small or labels are noisy. Padding choices can create boundary artifacts. CNNs encode locality, so they may struggle with long-range dependencies unless depth, dilation, attention, or other mechanisms provide a wider context.

\paragraph{Exercises.}
\begin{enumerate}[label=\textbf{40.1.\arabic*.}, leftmargin=*]
\item \textbf{Mechanical.} For \(x=(2,0,1,4)\) and \(K=(1,-1)\), compute the valid convolution \(y_t=x_t-x_{t+1}\).
\item \textbf{Proof/derivation.} Prove the one-dimensional translation-equivariance identity \(K\ast(T_s x)=T_s(K\ast x)\) away from boundaries.
\item \textbf{Numerical/computational.} A \(3\times3\) image convolution has \(C_{\mathrm{in}}=16\) and \(C_{\mathrm{out}}=32\). How many kernel weights are learned, excluding biases?
\item \textbf{Interpretation.} Explain why local convolutional filters are a reasonable first model for early visual processing but not a complete model of the visual cortex.
\item \textbf{Misconception/debugging.} A student says, ``CNNs recognize an object regardless of where it appears because convolution is invariant.'' Correct the statement using equivariance, pooling, and finite-boundary effects.
\end{enumerate}

\subsection{40.2 Recurrent neural networks and state-space models}

\paragraph{Guiding question.}
How can a neural network process a sequence by maintaining a hidden state that summarizes the past?

\paragraph{Beginner-facing intuition.}
Some data arrive in order: words in a sentence, spikes over time, voltage traces, movements, or sensory streams. A recurrent neural network (RNN) reads one input at a time and updates a hidden state. The hidden state is a memory: it carries information from previous inputs into the next computation.

This makes RNNs natural relatives of state-space models and dynamical systems. The state update can be learned, nonlinear, and driven by input. Training requires backpropagation through time, which repeatedly multiplies derivatives along the sequence and can produce vanishing or exploding gradients.

\paragraph{Formal definitions and notation.}
Let \(x_t\in\mathbb{R}^{d_x}\) be an input at time \(t\), \(h_t\in\mathbb{R}^{d_h}\) a hidden state, and \(o_t\in\mathbb{R}^{d_o}\) an output. A basic RNN is
\[
h_t=\phi(W_h h_{t-1}+W_x x_t+b_h),
\qquad
o_t=W_o h_t+b_o,
\]
where
\[
W_h\in\mathbb{R}^{d_h\times d_h},
\quad
W_x\in\mathbb{R}^{d_h\times d_x},
\quad
W_o\in\mathbb{R}^{d_o\times d_h}.
\]
A discrete state-space model has the more general form
\[
h_t=F_\theta(h_{t-1},x_t,\epsilon_t),
\qquad
y_t=G_\theta(h_t,\eta_t),
\]
with process noise \(\epsilon_t\) and observation noise \(\eta_t\) when stochasticity is included. Continuous-time versions use differential equations such as
\[
\frac{dh}{dt}=f_\theta(h(t),x(t),t).
\]

\paragraph{Core derivation: memory and gradient scale in a linear RNN.}
For the scalar recurrence
\[
h_t=ah_{t-1}+bx_t,\qquad h_0=0,
\]
repeated substitution gives
\[
h_t=b\sum_{k=1}^{t}a^{t-k}x_k.
\]
If \(|a|<1\), older inputs are exponentially forgotten. If \(|a|>1\), old inputs can be amplified and the state may explode. The same coefficient controls gradients:
\[
\frac{\partial h_t}{\partial h_k}=a^{t-k}.
\]
For vector RNNs, powers of the recurrent Jacobian play the same role. Eigenvalues or singular values below one cause vanishing signals; values above one cause exploding signals. Gated architectures such as LSTMs and GRUs introduce learned paths that help preserve or reset memory.

\paragraph{Concrete example.}
Let \(h_t=0.5h_{t-1}+x_t\), \(h_0=0\), and inputs \(x_1=2,x_2=0,x_3=4\). Then
\[
h_1=2,\qquad
h_2=1,\qquad
h_3=4.5.
\]
The first input contributes \(0.5^2\cdot2=0.5\) to \(h_3\), while the third input contributes \(4\). The hidden state is a weighted memory of the past.

\paragraph{Computational neuroscience example.}
Recurrent models are used to describe persistent activity, working memory, decision integration, and motor preparation. For example, a low-dimensional hidden state may represent a neural population trajectory during a delayed-response task. The RNN state \(h_t\) is a modeling variable; it may summarize recorded neural activity without being identical to any single biological state.

\paragraph{AI/ML example.}
RNNs and gated RNNs are used for time-series forecasting, speech, event streams, and sequence labeling. In language modeling they have largely been replaced by transformers for many tasks, but the state-space viewpoint remains important in control, neural data modeling, and newer sequence models that combine recurrence with efficient long-context computation.

\paragraph{Common misconceptions and failure modes.}
The hidden state is not guaranteed to remember everything; it is a finite-dimensional summary trained for a task. Stability is not always good: a model that forgets too quickly cannot use long context. Instability is not always useful: exploding states make training and prediction unreliable. Backpropagation through time is not a new calculus rule; it is the chain rule applied to an unrolled computational graph.

\paragraph{Exercises.}
\begin{enumerate}[label=\textbf{40.2.\arabic*.}, leftmargin=*]
\item \textbf{Mechanical.} For \(h_t=0.8h_{t-1}+x_t\), \(h_0=1\), and \(x_1=0,x_2=2\), compute \(h_1\) and \(h_2\).
\item \textbf{Proof/derivation.} Derive \(h_t=b\sum_{k=1}^{t}a^{t-k}x_k\) for the scalar linear recurrence with \(h_0=0\).
\item \textbf{Numerical/computational.} If \(\partial h_t/\partial h_{t-1}=0.9\) at every step, what is the gradient multiplier across \(50\) steps? Repeat for \(1.05\).
\item \textbf{Interpretation.} In a neural working-memory model, what would \(|a|\) near \(1\) mean in the scalar recurrence analogy?
\item \textbf{Misconception/debugging.} A student says, ``An RNN has memory, so it can use arbitrarily old inputs.'' Explain the role of finite state, training, and gradient decay.
\end{enumerate}

\subsection{40.3 Attention and transformers}

\paragraph{Guiding question.}
How can a network let each token choose which other tokens to read from, rather than compressing the past into a single recurrent state?

\paragraph{Beginner-facing intuition.}
Attention is content-based routing. Each item in a sequence produces a query, a key, and a value. Queries ask what information is needed; keys advertise what information is available; values contain the information to be mixed. A token attends to other tokens by comparing its query with their keys and forming a weighted average of their values.

Transformers stack attention with feedforward layers, residual connections, normalization, and positional information. Unlike a basic RNN, a self-attention layer can connect any pair of tokens in one step, although the computational cost grows with the number of token pairs.

\paragraph{Formal definitions and notation.}
Let a sequence representation be
\[
X\in\mathbb{R}^{T\times d_{\mathrm{model}}},
\]
where \(T\) is token count and \(d_{\mathrm{model}}\) is feature dimension. For one attention head,
\[
Q=XW_Q,\qquad K=XW_K,\qquad V=XW_V,
\]
with
\[
Q,K\in\mathbb{R}^{T\times d_k},\qquad V\in\mathbb{R}^{T\times d_v}.
\]
Scaled dot-product attention is
\[
\boxed{
\operatorname{Attn}(X)
=A V,\qquad
A=\operatorname{softmax}_{\mathrm{row}}\left(\frac{QK^\top}{\sqrt{d_k}}\right).
}
\]
The matrix \(A\in\mathbb{R}^{T\times T}\) has nonnegative rows summing to \(1\). Multi-head attention performs several such operations in parallel and combines their outputs. Positional encodings or position-dependent biases are needed because attention alone treats the sequence as a set of token features with pairwise content scores.

\paragraph{Core derivation: attention rows are adaptive convex combinations.}
For token \(i\), define scores
\[
s_{ij}=\frac{q_i^\top k_j}{\sqrt{d_k}}.
\]
The row softmax gives
\[
A_{ij}=\frac{\exp(s_{ij})}{\sum_{m=1}^{T}\exp(s_{im})}.
\]
Therefore \(A_{ij}\geq0\) and
\[
\sum_{j=1}^{T}A_{ij}=1.
\]
The output at token \(i\) is
\[
y_i=\sum_{j=1}^{T}A_{ij}v_j.
\]
Thus each output is a convex combination of value vectors, but the weights depend on the current content through query-key similarities. The factor \(1/\sqrt{d_k}\) keeps dot products from becoming too large as feature dimension grows.

\paragraph{Concrete example.}
Let one query have scores \(s=(2,0)\) for two tokens. Then
\[
A_1=\left(\frac{e^2}{e^2+1},\frac{1}{e^2+1}\right)\approx(0.881,0.119).
\]
If values are \(v_1=(1,0)\) and \(v_2=(0,2)\), the attended output is
\[
y_1=0.881(1,0)+0.119(0,2)=(0.881,0.238).
\]
The token mostly reads from the first value but keeps a small contribution from the second.

\paragraph{Computational neuroscience example.}
Selective attention in neuroscience concerns task- and context-dependent modulation of sensory processing. Transformer attention is a precise differentiable routing mechanism. The terms are related by a broad idea of selective weighting, but a transformer attention matrix is not direct evidence for biological attention weights. A careful comparison asks which signals are queries, keys, and values in a specified model.

\paragraph{AI/ML example.}
In language models, attention lets a token use information from earlier tokens such as subject nouns, definitions, or code variables. In vision transformers, image patches are tokens. In neural data analysis, attention can model interactions among time bins, neurons, or trials, but the learned weights must be validated rather than interpreted automatically as causal influence.

\paragraph{Common misconceptions and failure modes.}
Attention weights are not always explanations; values, downstream layers, and residual paths also affect outputs. Without positional information, self-attention is permutation equivariant and cannot know token order. Attention has quadratic token-pair cost unless modified. Very sharp attention can ignore useful context, while diffuse attention can blur information.

\paragraph{Exercises.}
\begin{enumerate}[label=\textbf{40.3.\arabic*.}, leftmargin=*]
\item \textbf{Mechanical.} If \(X\in\mathbb{R}^{12\times64}\), \(W_Q,W_K\in\mathbb{R}^{64\times16}\), and \(W_V\in\mathbb{R}^{64\times32}\), state the shapes of \(Q\), \(K\), \(V\), \(QK^\top\), and \(AV\).
\item \textbf{Proof/derivation.} Prove that each row of the softmax attention matrix \(A\) sums to \(1\).
\item \textbf{Numerical/computational.} Compute the softmax weights for scores \((1,1,0)\), then compute the weighted average of scalar values \((2,4,10)\).
\item \textbf{Interpretation.} Why does a transformer need positional information for ordered language or time-series data?
\item \textbf{Misconception/debugging.} A researcher interprets a high attention weight from token \(i\) to token \(j\) as proof that token \(j\) caused the model's decision. Give two reasons this conclusion may fail.
\end{enumerate}

\subsection{40.4 Graph neural networks}

\paragraph{Guiding question.}
How can a neural network process data whose objects are nodes connected by edges rather than entries on a grid or positions in a sequence?

\paragraph{Beginner-facing intuition.}
Graphs represent relationships: neurons connected by synapses, brain regions connected by tracts, atoms connected by bonds, users connected by interactions, or variables connected by dependencies. A graph neural network (GNN) updates each node by combining its own features with messages from neighboring nodes. The same update rule is reused across nodes, so the network does not depend on arbitrary node ordering.

The central symmetry is permutation equivariance. If we relabel the nodes and apply the same GNN, the output should be relabeled in the same way. The graph structure matters, but the names \(1,2,\ldots,N\) should not.

\paragraph{Formal definitions and notation.}
Let \(G=(V,E)\) be a graph with \(N=|V|\) nodes. Node \(i\) has feature vector \(h_i^{(\ell)}\in\mathbb{R}^{d_\ell}\) at layer \(\ell\). A general message-passing layer is
\[
m_i^{(\ell)}=\operatorname{AGG}_{j\in\mathcal{N}(i)}
\phi_\ell\!\left(h_i^{(\ell)},h_j^{(\ell)},e_{ij}\right),
\]
\[
h_i^{(\ell+1)}
=\psi_\ell\!\left(h_i^{(\ell)},m_i^{(\ell)}\right),
\]
where \(\mathcal{N}(i)\) is the neighborhood of node \(i\), \(e_{ij}\) is an optional edge feature, \(\operatorname{AGG}\) is a permutation-invariant aggregation such as sum, mean, or max, and \(\phi_\ell,\psi_\ell\) are learned maps.

A common linear graph convolution form is
\[
H^{(\ell+1)}=\sigma(\widetilde A H^{(\ell)}W_\ell),
\]
where \(H^{(\ell)}\in\mathbb{R}^{N\times d_\ell}\), \(W_\ell\in\mathbb{R}^{d_\ell\times d_{\ell+1}}\), and \(\widetilde A\) is a normalized adjacency matrix including self-loops.

\paragraph{Core derivation: permutation equivariance.}
Let \(P\in\mathbb{R}^{N\times N}\) be a permutation matrix that relabels nodes. Under relabeling,
\[
H\mapsto PH,\qquad \widetilde A\mapsto P\widetilde A P^\top.
\]
For the linear graph convolution,
\[
(P\widetilde A P^\top)(PH)W
=P\widetilde A (P^\top P)HW
=P\widetilde A HW.
\]
Applying the activation row-wise gives
\[
\sigma((P\widetilde A P^\top)(PH)W)=P\sigma(\widetilde AHW).
\]
Thus relabeling before the layer is the same as applying the layer and then relabeling the output. This is permutation equivariance.

\paragraph{Concrete example.}
Consider a three-node path \(1-2-3\) with scalar features
\[
h_1=1,\qquad h_2=3,\qquad h_3=5.
\]
Use mean aggregation with self-loops:
\[
h_i'=\text{mean}\{h_j:j\in\mathcal{N}(i)\cup\{i\}\}.
\]
Then
\[
h_1'=(1+3)/2=2,\qquad
h_2'=(1+3+5)/3=3,\qquad
h_3'=(3+5)/2=4.
\]
The middle node averages both neighbors; endpoints average themselves and the middle node.

\paragraph{Computational neuroscience example.}
GNNs can model connectome data by placing brain regions, neurons, or recording sites on nodes and anatomical or functional connections on edges. A message-passing step resembles local interaction along a network. This is a modeling abstraction: real synapses have direction, delay, neurotransmitter type, plasticity, and dynamics that a simple undirected GNN may omit.

\paragraph{AI/ML example.}
In molecular property prediction, atoms are nodes, bonds are edges, and message passing builds atom representations that are pooled into a molecule representation. In recommendation systems, users and items form bipartite graphs. In program analysis, variables or syntax elements can be graph nodes. The same mathematics applies once node features, edge features, and aggregation rules are defined.

\paragraph{Common misconceptions and failure modes.}
Graph neural networks are not automatically more expressive than all graph algorithms; standard message passing has known limitations related to distinguishing certain graph structures. Repeated averaging can cause oversmoothing, where node representations become too similar. If labels depend on long-range or global graph properties, shallow local message passing may fail. If the graph edges are noisy or inferred, the GNN may learn artifacts of graph construction.

\paragraph{Exercises.}
\begin{enumerate}[label=\textbf{40.4.\arabic*.}, leftmargin=*]
\item \textbf{Mechanical.} For a graph with \(N=10\), node-feature dimension \(8\), and output dimension \(16\), state the shape of \(H\), \(W\), and \(\widetilde AHW\) in the graph convolution formula.
\item \textbf{Proof/derivation.} Prove the permutation-equivariance identity for \(H^{+}=\sigma(\widetilde AHW)\) using a permutation matrix \(P\).
\item \textbf{Numerical/computational.} For the three-node path example, repeat the mean-aggregation update starting from \(h=(0,6,0)\).
\item \textbf{Interpretation.} In a connectome GNN, what do nodes, edges, node features, and messages represent? State one biological detail the model may ignore.
\item \textbf{Misconception/debugging.} A student applies ten mean-aggregation GNN layers and finds all node embeddings nearly identical. Name the failure mode and explain why repeated local averaging can cause it.
\end{enumerate}

\paragraph{Closing bridge.}
Architectures define the differentiable maps that optimization trains. The next chapter changes the learning problem: instead of predicting labels or outputs, generative deep learning models distributions over data. CNNs, RNNs, transformers, and GNNs will appear there as encoders, decoders, generators, discriminators, and denoising networks.

\clearpage
